% Part IV-A — How Verification Works: Types and Effects (Slides 61–80)

% ---------------------------------------------------------------------------
% Slide 61: Section divider
% ---------------------------------------------------------------------------
\begin{frame}
\frametitle{}
\vfill
\begin{center}
  {\Huge\color{jugeoBlue}\textbf{Part IV}}\\[12pt]
  {\LARGE\color{jugeoDark}How Verification Works in JuGeo}\\[20pt]
  \begin{tikzpicture}
    \draw[jugeoBlue, line width=2pt] (-3,0) -- (3,0);
  \end{tikzpicture}\\[12pt]
  {\large\color{gray}\itshape Refinement types, effects, and the verification pipeline}
\end{center}
\vfill
\end{frame}

% ---------------------------------------------------------------------------
% Slide 62: Overview — The Verification Pipeline
% ---------------------------------------------------------------------------
\begin{frame}
\frametitle{Overview: The Verification Pipeline}
\begin{center}
\begin{tikzpicture}[
    box/.style={draw=jugeoBlue, rounded corners, fill=jugeoLight,
                minimum width=26mm, minimum height=12mm, font=\small\bfseries,
                line width=1pt, align=center, text=jugeoDark},
    inp/.style={draw=jugeoGreen, rounded corners, fill=jugeoGreen!10,
                minimum width=22mm, minimum height=10mm, font=\small\bfseries,
                line width=1pt, text=jugeoGreen},
    outp/.style={draw=jugeoBlue, rounded corners, fill=jugeoBlue!15,
                 minimum width=30mm, minimum height=12mm, font=\small\bfseries,
                 line width=1pt, align=center, text=jugeoDark},
    arr/.style={-{Stealth[length=3mm]}, thick, jugeoBlue}
  ]
  \node[inp]  (code) at (-4.2,0.8)  {Python Code};
  \node[inp]  (spec) at (-4.2,-0.8) {Specification};
  \node[box]  (parse) at (-1,0)     {Parse \&\\Coordinate};
  \node[box]  (refine) at (2,0)     {Refine \&\\Verify};
  \node[outp] (out)  at (5.5,0)     {Proof-Carrying\\Python};

  \draw[arr] (code.east)  -- (parse.west |- code);
  \draw[arr] (spec.east)  -- (parse.west |- spec);
  \draw[arr] (parse) -- (refine);
  \draw[arr] (refine) -- (out);
\end{tikzpicture}
\end{center}
\vfill
\begin{itemize}
  \item \jugmark{Input}: your Python code + a specification of desired behavior
  \item \jugmark{Pipeline}: parse, assign coordinates, refine types, verify
  \item \jugmark{Output}: original Python with machine-checked proofs attached
\end{itemize}
\end{frame}

% ---------------------------------------------------------------------------
% Slide 63: What Are Refinement Types?
% ---------------------------------------------------------------------------
\begin{frame}
\frametitle{What Are Refinement Types?}
\begin{block}{Idea}
  A \jugmark{refinement type} = a base type \textbf{+} a logical condition.
\end{block}
\vspace{8pt}
\begin{center}
\begin{tikzpicture}
  \node[draw=jugeoBlue, rounded corners, fill=jugeoLight, inner sep=10pt,
        font=\Large\ttfamily] {%
    \{x : int \,|\, x >= 0\}%
  };
\end{tikzpicture}
\end{center}
\vspace{6pt}
\begin{columns}[T]
\begin{column}{0.48\textwidth}
\begin{alertblock}{\texttt{int} alone}
  \texttt{x} could be $-1$, $0$, $42$, $\ldots$
\end{alertblock}
\end{column}
\begin{column}{0.48\textwidth}
\begin{exampleblock}{\texttt{\{x:int | x>=0\}}}
  Only $0, 1, 2, \ldots$ allowed
\end{exampleblock}
\end{column}
\end{columns}
\vspace{6pt}
\begin{center}
  \itshape More precise than plain types --- catches more bugs at compile time.
\end{center}
\end{frame}

% ---------------------------------------------------------------------------
% Slide 64: Refinement Types in F*
% ---------------------------------------------------------------------------
\begin{frame}[fragile]
\frametitle{Refinement Types in F*}
\begin{exampleblock}{F* syntax}
\begin{lstlisting}
val abs : x:int -> Tot (y:int{y >= 0})
let abs x = if x >= 0 then x else -x
\end{lstlisting}
\end{exampleblock}
\vspace{6pt}
\begin{itemize}
  \item F* pioneered refinement types with \jugmark{SMT (Z3)} discharge
  \item Write the refinement $\to$ Z3 checks it automatically
  \item Powerful --- but two catches:
\end{itemize}
\vspace{6pt}
\begin{columns}[T]
\begin{column}{0.48\textwidth}
\begin{alertblock}{Catch 1}
  You must write in \textbf{F*'s language}, not Python
\end{alertblock}
\end{column}
\begin{column}{0.48\textwidth}
\begin{alertblock}{Catch 2}
  Trust is \textbf{binary}: proved or not --- no partial evidence
\end{alertblock}
\end{column}
\end{columns}
\end{frame}

% ---------------------------------------------------------------------------
% Slide 65: Refinement Types in JuGeo
% ---------------------------------------------------------------------------
\begin{frame}
\frametitle{Refinement Types in JuGeo}
\begin{block}{Key Difference}
  Every refinement lives at a \jugmark{coordinate} with its own \jugmark{evidence bundle}.
\end{block}
\vspace{6pt}
\begin{center}
\begin{tikzpicture}[
    node distance=4pt,
    evbox/.style={draw=jugeoBlue, rounded corners=3pt, fill=jugeoLight,
                  minimum width=52mm, font=\footnotesize, inner sep=5pt}
  ]
  \node[font=\small\ttfamily\bfseries, text=jugeoDark]
    (coord) {module.function.arg\_x};
  \node[evbox, below=8pt of coord] (z3)
    {\textcolor{jugeoGreen}{\texttt{SOLVER\_DISCHARGED}} --- Z3 proved it};
  \node[evbox, below=4pt of z3] (rt)
    {\textcolor{jugeoOrange}{\texttt{RUNTIME\_WITNESSED}} --- tests confirm it};
  \node[evbox, below=4pt of rt] (ai)
    {\textcolor{jugeoBlue}{\texttt{COPILOT\_SUGGESTED}} --- AI thinks it holds};
  \node[draw=jugeoDark, dashed, rounded corners=6pt, fit=(z3)(rt)(ai),
        inner sep=5pt, label={[font=\scriptsize, text=jugeoDark]right:evidence bundle}] {};
\end{tikzpicture}
\end{center}
\vspace{4pt}
\begin{center}
  \itshape Three sources, tracked separately --- combine or compare.
\end{center}
\end{frame}

% ---------------------------------------------------------------------------
% Slide 66: The Refinement Pipeline
% ---------------------------------------------------------------------------
\begin{frame}
\frametitle{The Refinement Pipeline}
\begin{center}
\begin{tikzpicture}[
    stage/.style={draw=jugeoBlue, rounded corners, fill=jugeoLight,
                  minimum width=32mm, minimum height=10mm,
                  font=\footnotesize\bfseries, align=center, text=jugeoDark, line width=1pt},
    arr/.style={-{Stealth[length=3mm]}, thick, jugeoBlue}
  ]
  \node[stage] (rsb)  at (0,0)    {RefinementSort\\Builder};
  \node[stage] (pn)   at (3.8,0)  {Predicate\\Normalizer};
  \node[stage] (cl)   at (7.6,0)  {Constraint\\Lifter};
  \node[stage] (rte)  at (11.4,0) {RefinementType\\Encoder};

  \draw[arr] (rsb) -- (pn);
  \draw[arr] (pn)  -- (cl);
  \draw[arr] (cl)  -- (rte);

  \node[below=6pt of rsb, font=\scriptsize, text=gray, align=center]
    {Assigns sorts to\\each coordinate};
  \node[below=6pt of pn, font=\scriptsize, text=gray, align=center]
    {Standardizes\\predicates};
  \node[below=6pt of cl, font=\scriptsize, text=gray, align=center]
    {Lifts constraints\\across scopes};
  \node[below=6pt of rte, font=\scriptsize, text=gray, align=center]
    {Emits SMT-LIB\\formulas};
\end{tikzpicture}
\end{center}
\vfill
\begin{enumerate}
  \item \jugmark{Sort}: determine the logical sort of each refinement
  \item \jugmark{Normalize}: rewrite predicates into canonical form
  \item \jugmark{Lift}: propagate constraints from callees to callers
  \item \jugmark{Encode}: translate into SMT-ready assertions for Z3
\end{enumerate}
\end{frame}

% ---------------------------------------------------------------------------
% Slide 67: What Are Effects?
% ---------------------------------------------------------------------------
\begin{frame}
\frametitle{What Are Effects?}
\begin{block}{Definition}
  An \jugmark{effect} is anything a function does \textbf{besides} returning a value.
\end{block}
\vspace{8pt}
\begin{center}
\begin{tikzpicture}[
    eff/.style={draw=jugeoBlue, rounded corners, fill=jugeoLight,
                minimum width=24mm, minimum height=10mm,
                font=\small\bfseries, text=jugeoDark}
  ]
  \node[eff] (exc)  at (0,0)     {Exceptions};
  \node[eff] (mut)  at (3.2,0)   {Mutation};
  \node[eff] (io)   at (6.4,0)   {I/O};
  \node[eff] (asyn) at (9.6,0)   {Async};
  \node[eff] (gen)  at (1.6,-1.4){Generators};
  \node[eff] (ctx)  at (4.8,-1.4){Context Mgrs};
  \node[eff] (meta) at (8.0,-1.4){Metaclasses};
\end{tikzpicture}
\end{center}
\vfill
\begin{center}
  \large Python is \jugmark{effect-rich} --- almost every function has side effects.\\[4pt]
  \normalsize Verification must account for \textbf{all} of them.
\end{center}
\end{frame}

% ---------------------------------------------------------------------------
% Slide 68: Why Effects Matter for Verification
% ---------------------------------------------------------------------------
\begin{frame}
\frametitle{Why Effects Matter for Verification}
\begin{center}
\begin{tikzpicture}[
    box/.style={draw=jugeoBlue, rounded corners, fill=jugeoLight,
                minimum width=28mm, minimum height=14mm,
                font=\small, align=center, text=jugeoDark, line width=1pt},
    arr/.style={-{Stealth[length=3mm]}, thick}
  ]
  \node[box] (fn) at (0,0) {\textbf{f(x)}\\call function};
  \node[box, fill=jugeoGreen!15] (ok) at (4,1) {\textbf{return y}\\normal path};
  \node[box, fill=jugeoRed!15]   (ex) at (4,-1) {\textbf{raise E}\\exception path};

  \draw[arr, jugeoGreen] (fn.east) -- (ok.west) node[midway, above, font=\scriptsize] {success};
  \draw[arr, jugeoRed]   (fn.east) -- (ex.west) node[midway, below, font=\scriptsize] {failure};
\end{tikzpicture}
\end{center}
\vfill
\begin{itemize}
  \item A function that can throw must be verified along \jugmark{both} paths
  \item Same for \texttt{async} (suspend/resume), generators (\texttt{yield}/\texttt{send}),
        context managers (\texttt{\_\_enter\_\_}/\texttt{\_\_exit\_\_})
  \item \textbf{Ignoring effects = ignoring half the behavior}
\end{itemize}
\end{frame}

% ---------------------------------------------------------------------------
% Slide 69: Effects in F* — Dijkstra Monads
% ---------------------------------------------------------------------------
\begin{frame}[fragile]
\frametitle{Effects in F*: Dijkstra Monads}
\begin{block}{F*'s Approach}
  Model each effect as a \jugmark{Dijkstra monad} --- a mathematical object
  encoding the effect's weakest-precondition transformer.
\end{block}
\vspace{8pt}
\begin{exampleblock}{Example: F* stateful effect}
\begin{lstlisting}
effect ST (a:Type) (pre: heap -> Type)
         (post: heap -> a -> heap -> Type) =
  STATE a (fun p h -> pre h /\ forall x h'. post h x h' ==> p x h')
\end{lstlisting}
\end{exampleblock}
\vspace{4pt}
\begin{columns}[T]
\begin{column}{0.48\textwidth}
\begin{exampleblock}{Strengths}
  \begin{itemize}
    \item Mathematically precise
    \item Compositional
  \end{itemize}
\end{exampleblock}
\end{column}
\begin{column}{0.48\textwidth}
\begin{alertblock}{Limitations}
  \begin{itemize}
    \item Very abstract notation
    \item No native Python effects
  \end{itemize}
\end{alertblock}
\end{column}
\end{columns}
\end{frame}

% ---------------------------------------------------------------------------
% Slide 70: Effects in JuGeo — Geometric, Not Monadic
% ---------------------------------------------------------------------------
\begin{frame}
\frametitle{Effects in JuGeo: Geometric, Not Monadic}
\begin{columns}[T]
\begin{column}{0.48\textwidth}
\begin{alertblock}{Monadic (F*)}
  \begin{itemize}
    \item Effects wrapped in monads
    \item Compose via \textbf{monad transformers}
    \item Encoding overhead
  \end{itemize}
\end{alertblock}
\end{column}
\begin{column}{0.48\textwidth}
\begin{exampleblock}{Geometric (JuGeo)}
  \begin{itemize}
    \item Effects are \jugmark{sections} in the site
    \item Compose via \jugmark{gluing}
    \item Native Python semantics
  \end{itemize}
\end{exampleblock}
\end{column}
\end{columns}
\vspace{10pt}
\begin{center}
\begin{tikzpicture}
  \node[fill=jugeoLight, rounded corners, inner sep=8pt, text width=11cm, align=center,
        font=\itshape]
    {Key insight: effects compose via \jugmark{gluing}, not monad transformers.\\[2pt]
     Each effect is a geometric structure on the coordinate space.};
\end{tikzpicture}
\end{center}
\end{frame}

% ---------------------------------------------------------------------------
% Slide 71: Python Exceptions as Coordinate Forks
% ---------------------------------------------------------------------------
\begin{frame}
\frametitle{Python Exceptions as Coordinate Forks}
\begin{center}
\begin{tikzpicture}[
    node distance=8mm,
    cbox/.style={draw=jugeoBlue, rounded corners, fill=jugeoLight,
                 minimum width=30mm, minimum height=10mm,
                 font=\small\bfseries, align=center, text=jugeoDark, line width=1pt},
    arr/.style={-{Stealth[length=3mm]}, thick}
  ]
  \node[cbox] (try) {try block};
  \node[cbox, fill=jugeoGreen!15, below right=10mm and 12mm of try] (ok) {normal section};
  \node[cbox, fill=jugeoRed!15,   below left=10mm and 12mm of try]  (ex) {except section};
  \node[cbox, below=30mm of try] (join) {join point};

  \draw[arr, jugeoGreen] (try.south east) -- (ok.north)
    node[midway, right, font=\scriptsize] {no exception};
  \draw[arr, jugeoRed]   (try.south west) -- (ex.north)
    node[midway, left, font=\scriptsize] {exception raised};
  \draw[arr, jugeoGreen] (ok.south) -- (join.north east);
  \draw[arr, jugeoRed]   (ex.south) -- (join.north west);

  \node[right=4mm of join, font=\scriptsize\itshape, text=jugeoDark]
    {sections glued here};
\end{tikzpicture}
\end{center}
\vspace{2pt}
\begin{center}
  \texttt{try/except} creates a \jugmark{fork} in coordinate space ---\\
  two branches that must be \jugmark{glued} at the join point.
\end{center}
\end{frame}

% ---------------------------------------------------------------------------
% Slide 72: Exception Example
% ---------------------------------------------------------------------------
\begin{frame}[fragile]
\frametitle{Exception Example}
\begin{columns}[T]
\begin{column}{0.48\textwidth}
\begin{exampleblock}{Python code}
\begin{lstlisting}
try:
    result = 10 / x
except ZeroDivisionError:
    result = 0
\end{lstlisting}
\end{exampleblock}
\end{column}
\begin{column}{0.48\textwidth}
\begin{block}{JuGeo model}
\begin{tikzpicture}[
    br/.style={draw=jugeoBlue, rounded corners=3pt, fill=jugeoLight,
               minimum width=38mm, font=\scriptsize, inner sep=4pt,
               align=left, text=jugeoDark}
  ]
  \node[br, fill=jugeoGreen!12] (b1) at (0,0)
    {\texttt{x != 0}\\[1pt]\texttt{result = 10/x}};
  \node[br, fill=jugeoRed!12, below=4pt of b1] (b2)
    {\texttt{x == 0}\\[1pt]\texttt{result = 0}};
  \node[below=4pt of b2, font=\scriptsize\itshape, text=jugeoDark]
    {Both branches verified};
\end{tikzpicture}
\end{block}
\end{column}
\end{columns}
\vfill
\begin{itemize}
  \item \jugmark{Branch 1}: $x \neq 0$ --- normal division, \texttt{result = 10/x}
  \item \jugmark{Branch 2}: $x = 0$ --- exception caught, \texttt{result = 0}
  \item Both sections glued $\Rightarrow$ \texttt{result} always has a valid value
\end{itemize}
\end{frame}

% ---------------------------------------------------------------------------
% Slide 73: Python Async as Suspended Morphisms
% ---------------------------------------------------------------------------
\begin{frame}
\frametitle{Python Async as Suspended Morphisms}
\begin{center}
\begin{tikzpicture}[
    seg/.style={draw=jugeoBlue, rounded corners, fill=jugeoLight,
                minimum width=22mm, minimum height=9mm,
                font=\footnotesize\bfseries, text=jugeoDark, line width=1pt},
    susp/.style={draw=jugeoOrange, dashed, thick, rounded corners,
                 minimum width=12mm, minimum height=9mm,
                 font=\footnotesize\bfseries, text=jugeoOrange},
    arr/.style={-{Stealth[length=3mm]}, thick, jugeoBlue}
  ]
  \node[seg]  (s1)  at (0,0)      {compute A};
  \node[susp] (aw1) at (3,0)      {await};
  \node[seg]  (s2)  at (6,0)      {compute B};
  \node[susp] (aw2) at (9,0)      {await};
  \node[seg]  (s3)  at (12,0)     {compute C};

  \draw[arr] (s1) -- (aw1);
  \draw[arr, jugeoOrange, dashed] (aw1) -- (s2);
  \draw[arr] (s2) -- (aw2);
  \draw[arr, jugeoOrange, dashed] (aw2) -- (s3);
\end{tikzpicture}
\end{center}
\vfill
\begin{description}[\textbf{Cancellation}]
  \item[\textbf{Coroutine}] A \jugmark{suspended section morphism} ---
    paused mid-execution
  \item[\textbf{await}] A resume point where control may transfer
  \item[\textbf{Scheduling}] Morphism composition across tasks
  \item[\textbf{Cancellation}] An \textcolor{jugeoRed}{obstruction} ---
    the section cannot be completed
\end{description}
\end{frame}

% ---------------------------------------------------------------------------
% Slide 74: Python Generators as Fiber Sequences
% ---------------------------------------------------------------------------
\begin{frame}[fragile]
\frametitle{Python Generators as Fiber Sequences}
\begin{columns}[T]
\begin{column}{0.45\textwidth}
\begin{exampleblock}{Generator code}
\begin{lstlisting}
def counter(n):
    i = 0
    while i < n:
        yield i
        i += 1
\end{lstlisting}
\end{exampleblock}
\end{column}
\begin{column}{0.52\textwidth}
\begin{block}{JuGeo model}
\begin{tikzpicture}[
    ynode/.style={circle, draw=jugeoBlue, fill=jugeoLight,
                  minimum size=7mm, font=\scriptsize\bfseries,
                  text=jugeoDark},
    arr/.style={-{Stealth[length=2.5mm]}, thick, jugeoBlue}
  ]
  \node[ynode] (y0) at (0,0)    {$y_0$};
  \node[ynode] (y1) at (1.5,0)  {$y_1$};
  \node[ynode] (y2) at (3,0)    {$y_2$};
  \node[font=\large, text=jugeoBlue] at (4.2,0) {$\cdots$};

  \draw[arr] (y0) -- (y1);
  \draw[arr] (y1) -- (y2);

  \node[below=5pt of y0, font=\scriptsize, text=gray] {yield 0};
  \node[below=5pt of y1, font=\scriptsize, text=gray] {yield 1};
  \node[below=5pt of y2, font=\scriptsize, text=gray] {yield 2};
\end{tikzpicture}
\end{block}
\end{column}
\end{columns}
\vfill
\begin{itemize}
  \item \texttt{yield} emits a \jugmark{partial section} --- one element of the fiber
  \item \texttt{send()} injects a value into the generator fiber
  \item \texttt{close()} terminates the section
  \item Each yield point is a \jugmark{coordinate} in the sequence
\end{itemize}
\end{frame}

% ---------------------------------------------------------------------------
% Slide 75: Python Context Managers as Covers
% ---------------------------------------------------------------------------
\begin{frame}
\frametitle{Python Context Managers as Covers}
\begin{center}
\begin{tikzpicture}[
    phase/.style={draw=jugeoBlue, rounded corners, fill=jugeoLight,
                  minimum width=28mm, minimum height=12mm,
                  font=\small\bfseries, align=center, text=jugeoDark, line width=1pt},
    arr/.style={-{Stealth[length=3mm]}, thick, jugeoBlue}
  ]
  \node[phase, fill=jugeoGreen!15] (enter) at (0,0)   {\_\_enter\_\_\\[-1pt]\scriptsize\normalfont setup};
  \node[phase]                      (body)  at (4.2,0) {body\\[-1pt]\scriptsize\normalfont execution};
  \node[phase, fill=jugeoOrange!15] (exit)  at (8.4,0) {\_\_exit\_\_\\[-1pt]\scriptsize\normalfont teardown};

  \draw[arr] (enter) -- (body);
  \draw[arr] (body) -- (exit);

  \node[draw=jugeoDark, dashed, rounded corners=8pt, fit=(enter)(body)(exit),
        inner sep=8pt,
        label={[font=\footnotesize\itshape, text=jugeoDark]above:covering family}] {};
\end{tikzpicture}
\end{center}
\vfill
\begin{itemize}
  \item A \texttt{with}-block is a \jugmark{covering family} with three phases
  \item \texttt{\_\_enter\_\_}: acquire the resource (setup morphism)
  \item \texttt{body}: use the resource (section)
  \item \texttt{\_\_exit\_\_}: release the resource (teardown morphism)
  \item The cover \jugmark{guarantees} teardown runs --- even on exception
\end{itemize}
\end{frame}

% ---------------------------------------------------------------------------
% Slide 76: Context Manager Example
% ---------------------------------------------------------------------------
\begin{frame}[fragile]
\frametitle{Context Manager Example}
\begin{columns}[T]
\begin{column}{0.45\textwidth}
\begin{exampleblock}{Python code}
\begin{lstlisting}
with open("file.txt") as f:
    data = f.read()
# f is closed here
\end{lstlisting}
\end{exampleblock}
\end{column}
\begin{column}{0.52\textwidth}
\begin{block}{JuGeo model}
\begin{tikzpicture}[
    ph/.style={draw=jugeoBlue, rounded corners=3pt, fill=jugeoLight,
               minimum width=40mm, font=\scriptsize, inner sep=4pt,
               align=left, text=jugeoDark}
  ]
  \node[ph, fill=jugeoGreen!12] (e) at (0,0)
    {\textbf{enter}: \texttt{f = open(\ldots)}};
  \node[ph, below=3pt of e] (b)
    {\textbf{body}: \texttt{data = f.read()}};
  \node[ph, fill=jugeoOrange!12, below=3pt of b] (x)
    {\textbf{exit}: \texttt{f.close()}};
\end{tikzpicture}
\end{block}
\end{column}
\end{columns}
\vfill
\begin{block}{Temporal obligation}
  JuGeo encodes: \texttt{\_\_exit\_\_} \jugmark{MUST} execute after the body completes ---
  whether normally or via exception. This is a provable property, not just convention.
\end{block}
\end{frame}

% ---------------------------------------------------------------------------
% Slide 77: Python Metaclasses
% ---------------------------------------------------------------------------
\begin{frame}
\frametitle{Python Metaclasses}
\begin{center}
\begin{tikzpicture}[
    phase/.style={draw=jugeoBlue, rounded corners, fill=jugeoLight,
                  minimum width=30mm, minimum height=14mm,
                  font=\footnotesize\bfseries, align=center,
                  text=jugeoDark, line width=1pt},
    arr/.style={-{Stealth[length=3mm]}, thick, jugeoBlue}
  ]
  \node[phase, fill=jugeoGreen!15] (prep) at (0,0)
    {type.\_\_prepare\_\_()\\[-1pt]\scriptsize\normalfont create namespace};
  \node[phase] (body) at (4.5,0)
    {class body\\[-1pt]\scriptsize\normalfont populate namespace};
  \node[phase, fill=jugeoOrange!15] (new) at (9,0)
    {type.\_\_new\_\_()\\[-1pt]\scriptsize\normalfont create class object};

  \draw[arr] (prep) -- (body);
  \draw[arr] (body) -- (new);
\end{tikzpicture}
\end{center}
\vfill
\begin{itemize}
  \item Metaclass creation is a \jugmark{three-phase morphism sequence}
  \item Each phase transforms the namespace at a distinct coordinate
  \item JuGeo verifies the contract across all three phases
\end{itemize}
\vspace{4pt}
\begin{center}
\begin{tikzpicture}
  \node[fill=jugeoLight, rounded corners, inner sep=6pt, font=\footnotesize\itshape]
    {No other proof assistant models Python metaclass creation};
\end{tikzpicture}
\end{center}
\end{frame}

% ---------------------------------------------------------------------------
% Slide 78: Effect Comparison Table
% ---------------------------------------------------------------------------
\begin{frame}
\frametitle{Effect Comparison: JuGeo vs.\ the Big Three}
\begin{center}
\small
\renewcommand{\arraystretch}{1.3}
\begin{tabular}{l c c c c}
  \toprule
  \textbf{Effect} & \textbf{F*} & \textbf{LEAN} & \textbf{Coq} & \textbf{JuGeo} \\
  \midrule
  Exceptions        & custom & custom & --- &
    \textcolor{jugeoGreen}{\textbf{native}} \\
  State mutation     & Dijkstra & custom & --- &
    \textcolor{jugeoGreen}{\textbf{native}} \\
  Async / await     & --- & --- & --- &
    \textcolor{jugeoGreen}{\textbf{native}} \\
  Generators        & --- & --- & --- &
    \textcolor{jugeoGreen}{\textbf{native}} \\
  Context managers   & --- & --- & --- &
    \textcolor{jugeoGreen}{\textbf{native}} \\
  Metaclasses       & --- & --- & --- &
    \textcolor{jugeoGreen}{\textbf{native}} \\
  \bottomrule
\end{tabular}
\end{center}
\vfill
\begin{center}
  \jugmark{---} = no support \quad
  \textbf{custom} = requires manual encoding\\[4pt]
  \textcolor{jugeoGreen}{\textbf{native}} = built into the verification framework
\end{center}
\end{frame}

% ---------------------------------------------------------------------------
% Slide 79: Why Geometric Effects Are Better
% ---------------------------------------------------------------------------
\begin{frame}
\frametitle{Why Geometric Effects Are Better}
\begin{enumerate}
  \item \jugmark{Visible overlap conditions}
  \begin{itemize}
    \item Effects are explicit in the coordinate structure
    \item Not hidden inside monads --- you can \emph{see} where effects interact
  \end{itemize}
  \vspace{6pt}
  \item \jugmark{Native Python effects}
  \begin{itemize}
    \item No encoding step --- \texttt{try/except}, \texttt{async/await},
          \texttt{yield}, \texttt{with} all modeled directly
    \item Python semantics in, Python semantics out
  \end{itemize}
  \vspace{6pt}
  \item \jugmark{Persistent obstructions}
  \begin{itemize}
    \item When an effect \emph{fails} to compose, the obstruction is recorded
    \item Like a ``proof debt'' --- visible, trackable, resolvable
  \end{itemize}
\end{enumerate}
\end{frame}

% ---------------------------------------------------------------------------
% Slide 80: Effects Summary
% ---------------------------------------------------------------------------
\begin{frame}
\frametitle{Effects Summary}
\vfill
\begin{center}
\begin{tikzpicture}
  \node[draw=jugeoBlue, rounded corners=8pt, fill=jugeoLight,
        inner sep=14pt, text width=12cm, align=center,
        line width=1.5pt, font=\large] {%
    JuGeo is the \jugmark{first proof assistant} that natively models\\[6pt]
    \textbf{all} of Python's effect system:\\[10pt]
    \textcolor{jugeoGreen}{exceptions} \quad
    \textcolor{jugeoGreen}{async/await} \quad
    \textcolor{jugeoGreen}{generators}\\[6pt]
    \textcolor{jugeoGreen}{context managers} \quad
    \textcolor{jugeoGreen}{metaclasses}\\[10pt]
    --- as \jugmark{geometric structures} on the semantic site.%
  };
\end{tikzpicture}
\end{center}
\vfill
\begin{center}
  \itshape Next: how JuGeo's coordinate system enables incremental verification.
\end{center}
\end{frame}
